% Generated draft: requires mathematical and source review.
\documentclass[11pt]{article}
\usepackage[margin=25mm]{geometry}
\usepackage{fontspec}
\setmainfont{Latin Modern Roman}
\usepackage{amsmath,amssymb,mathtools,unicode-math}
\setmathfont{Latin Modern Math}
\usepackage{enumitem,needspace,xurl,hyperref}
\setlength{\parindent}{0pt}
\setlength{\parskip}{5pt}
\setlength{\emergencystretch}{3em}
\newcommand{\N}{\mathbb{N}}
\newcommand{\Z}{\mathbb{Z}}
\newcommand{\Q}{\mathbb{Q}}
\newcommand{\R}{\mathbb{R}}
\newcommand{\C}{\mathbb{C}}
\newcommand{\F}{\mathbb{F}}
\newcommand{\A}{\mathbb{A}}
\newcommand{\PP}{\mathbb{P}}
\newcommand{\E}{\mathbb{E}}
\newcommand{\eps}{\varepsilon}
\newcommand{\dd}{\,\mathrm{d}}
\begin{document}
\section*{UnsolvedMath ID: 30005579}
\textbf{Source problem code:} catalogue entry
\subsection{Definitions and mathematical statement}
Let \(M\) denote a smooth real four-dimensional manifold. A symplectic structure on \(M\) is a smooth differential two-form \(\omega\) satisfying \(\dd\omega=0\) and such that, for every \(p\in M\), the condition \(\omega_p(v,w)=0\) for all \(w\in T_pM\) implies \(v=0\). A Riemannian metric \(g\) on \(M\) is Einstein if there is a constant \(\lambda\in\R\) such that
\[
\operatorname{Ric}_g=\lambda g.
\]
\textbf{Exploratory problem.} For which additional smooth four-manifolds \(M\) can a symplectic structure \(\omega\) be used to obtain information about the existence or structure of Einstein metrics on \(M\)?

The supplied record does not identify the previously discussed examples relative to which manifolds are called ``additional,'' or specify what information about Einstein metrics is sought. It also does not specify compactness, connectedness, boundary conditions, completeness, the sign of \(\lambda\), or any compatibility between \(\omega\) and \(g\). These qualifications remain unresolved; no such restrictions or compatibility requirements are imposed by this draft. The record therefore supplies an exploratory question rather than a fully specified conjecture with a definite quantified conclusion.
\subsection{Short English statement}
Determine further four-manifolds on which symplectic geometry yields information about Einstein metrics. The comparison examples, manifold hypotheses, and intended conclusions are unspecified, so a more precise conjecture requires additional source context.
\subsection{Sources}
\begin{itemize}
\item Local UnsolvedMath registry, \texttt{data/unsolvedmath/problems.json}, permanent ID \(30005579\), title ``Symplectic Structures and Einstein Metrics,'' source problem code \texttt{catalogue entry}; supplied collection: Oberwolfach Reports Open Problems. No individual proposer is named.
\item Supplied source citation: \emph{Differentialgeometrie im Grossen (2024)}, Oberwolfach Reports, DOI \texttt{10.4171/owr/2023/29}, pp.\ 1617--1670. Supplied source URL and literature-source link, labeled ``original OWR report'': \url{https://doi.org/10.4171/owr/2023/29}. The record describes Oberwolfach Reports as documenting mathematical workshops and their open-problem discussions.
\item Supplied original wording: ``Are there other 4-manifolds where a symplectic structure helps reveal something about Einstein metrics?'' The record states that the report passage is retained verbatim as relevant context, but supplies no further surrounding mathematical discussion or precise passage location within the cited page range.
\item Literature assessment dated \texttt{2026-08-22}: the supplied curation labels the problem open and describes it as a broad exploratory question for which no general resolving theorem was identified. This is supplied dated curation derived from the OWR material, not an independently verified literature assessment; time-sensitive status and cited claims require independent verification.
\item The registry labels the statement \texttt{corrected\_verified} and reports that the official report was checked directly and that the cleaned statement preserves the complete mathematical target while restoring definitions and removing extraction artifacts or source-dependent wording. These are supplied verification claims, not checks performed for this draft. The proposed year and year tag are \(2023\), while the supplied citation names \(2024\); this discrepancy remains unresolved.
\end{itemize}
\end{document}
